Celem niniejszej pracy było opracowanie biblioteki
umożliwiającej automatyczną optymalizację
architektury w~pełni połączonych
wielowarstwowych percepronów
z~wykorzystaniem algorytmów genetycznych.
Projektowanie struktury sieci neuronowych
stanowi jeden z~kluczowych etapów procesu
budowy modeli uczenia maszynowego,
mający bezpośredni wpływ na ich skuteczność,
zdolność generalizacji oraz czas uczenia.
Dobór jej parametrów -- takich jak
liczba warstw czy liczba neuronów w~nich
-- jest jednakże najczęściej
realizowany metodą prób i~błędów,
co wymaga znacznego nakładu czasu
oraz wiedzy eksperckiej.

W~części teoretycznej przedstawiono zagadnienia
związane ze sztucznymi sieciami neuronowymi,
ze szczególnym uwzględnieniem architektury \acrshort{MLP}.
Omówiono również podstawy algorytmów ewolucyjnych,
ich mechanizmy działania oraz zastosowanie
algorytmów genetycznych w~zadaniach optymalizacyjnych.
A~następnie przeanalizowano istniejące
rozwiązania wykorzystywane do automatycznego
wyszukiwania architektur sieci neuronowych
oraz wskazano ich ograniczenia.

W~części praktycznej zostały przeanalizowane
wymagania związane z~projektem,
a~następnie wedle nich zaprojektowano
i~opisano bibliotekę \acrshort{GAAML},
zaimplementowaną w~języku Python,
pozwalającą na automatyczne dobieranie
zadawalającej struktury \acrshort{MLP}.
Opracowane rozwiązanie umożliwia
nienadzorowane optymalizowanie architektury
przy podaniu jedynie danych treningowych,
testowych i~ilości parametrów wejściowych,
pozwala natomiast również na podanie danych
walidacyjnych oraz zmienienie domyślnych
parametrów przeszukiwania przestrzeni
rozwiązań włącznie z~częścią
korygującą funkcji przystosowania.
Biblioteka została zaprojektowana
w~sposób umożliwiający jak najłatwiejszą:
możliwość przyszłego rozwijania
oraz integrację do~projektów.

W~celu oceny skuteczności opracowanego
rozwiązania przeprowadzono serię
eksperymentów na wybranych zbiorze danych.
Uzyskane wyniki potwierdziły możliwość efektywnego
wykorzystania algorytmów genetycznych do przeszukiwania
rozległej przestrzeni potencjalnych
architektur sieci neuronowych.
Przeprowadzone badania wykazały,
że bibloteka znajduje rozwiązania
lepsze niż czysto stochastyczna metoda,
więc potwierdzają zasadność
użycia algorytmów genetycznych
do projektowania sieci neuronowych.
